% !PRINTED 372
\weberpage{372}
% Source-faithful three-stroke Cayley configuration sign.
\providecommand{\weberthreeup}{\mathord{\tikz[baseline=-0.2ex,x=0.23em,y=0.23em,line width=0.45pt]{\draw (0,0)--(-1,1) (0,0)--(0,1.08) (0,0)--(1,1);}}}
If we consider the syzygetic complex triple
\weq{8}{
\begin{aligned}
&y_3 z_4,\ y_4 z_3,\ \cdot\ \cdot\ \cdot\ \cdot\\
&y_3 y_4,\ z_3 z_4,\ x_3 x_4\ \cdot\ \cdot\\
&y_3 z_3,\ y_4 z_4,\ \cdot\ \cdot\ \cdot\ \cdot,
\end{aligned}
}
in which all double tangents must occur, we find $[5\,8]$ and
$[5\,6]$ not in the last two complexes of (8),
because
\weq{9}{
\begin{aligned}
&y_3 z_3 [5\,8] = [1\,3]\,[2\,3]\,[5\,8],\quad y_3 y_4 [5\,8] = [1\,3]\,[1\,4]\,[5\,8]\\
&y_3 z_3 [5\,6] = [1\,3]\,[2\,3]\,[5\,6],\quad y_3 y_4 [5\,6] = [1\,3]\,[1\,4]\,[5\,6]
\end{aligned}
}
have the symbol $\vee\mid$ and are therefore azygetic. Consequently
$[5\,8]$ and $[5\,6]$ occur in the first of the complexes (8), and the
two triples (7) are syzygetic.

Finally we have also to consider the symbol $\sqcap$, which
again gives three types:
\weq{10}{
\begin{aligned}
&[1\,2]\,[2\,3]\,[3\,4] = y_2 z_3 [3\,4]\\
&[1\,2]\,[2\,3]\,[3\,8] = y_2 z_3 x_3\\
&[1\,8]\,[8\,2]\,[2\,3] = x_1 x_2 z_3.
\end{aligned}
}
These triples too are syzygetic, which for the last two is to be
seen immediately from the complex triple (3), and for
the first follows from the syzygetic complex triple
\[
\begin{aligned}
&y_2 z_3,\ x_1 x_3,\ \cdot\ \cdot\ \cdot\ \cdot\\
&y_2 x_1,\ z_3 x_3,\ \cdot\ \cdot\ \cdot\ \cdot\\
&y_2 x_3,\ z_3 x_1,\ \cdot\ \cdot\ \cdot\ \cdot,
\end{aligned}
\]
of which the last two do not contain $[3\,4]$, because
$y_2 x_1 [3\,4]$ and $y_2 x_3 [3\,4]$ both have the symbol $\vee\mid$.

Here, however, the exceptional position of the figure 8 has entirely
disappeared, and we arrive at the result:

\gesperrt{Among the triples of double tangents, those with the
symbols}
\[
\mid\mid\mid,\ \sqcap
\]
\gesperrt{are syzygetic, and those with the symbols}
\[
\weberthreeup,\ \triangle,\ \vee\mid
\]
\gesperrt{azygetic.}

From this one very easily obtains the symbols for all the
syzygetic and azygetic quadruples:
%CONT>>
